\chapter{Polygons and Three-Dimensional Representations}

The foundation of three-dimensional graphics lies in the effective representation and manipulation of polygonal geometry. This chapter provides a comprehensive exploration of how the Caffeine Engine handles polygon meshes, primitive generation, vertex data structures, and the mathematical frameworks that enable efficient rendering of complex three-dimensional objects.

\needspace{6\baselineskip}
\section{Polygon Mesh Fundamentals}

A polygon mesh is a collection of vertices, edges, and faces that together define the surface of a three-dimensional object. In real-time graphics engines like Caffeine, meshes are typically represented as collections of triangles (triangulated meshes) because triangles offer several advantages: they are planar, can be rasterized efficiently by modern GPUs, and can be decomposed from arbitrary polygons through triangulation algorithms.

\needspace{5\baselineskip}
\subsection{Vertex Attributes and Topology}

Each vertex in a mesh carries multiple attributes that influence how the mesh is rendered. At minimum, a vertex must contain a position in three-dimensional space. Additional attributes such as surface normals, texture coordinates, tangent vectors, and vertex colors provide the renderer with the necessary information to apply lighting, texturing, and other visual effects.

\begin{lstlisting}[language=C++, caption=Standard Vertex Structure]
struct Vertex {
    Vec3 position;           // Local coordinate (model space)
    Vec3 normal;             // Surface normal for lighting
    Vec2 texCoord;           // Texture sampling coordinate
    Vec3 tangent;            // Tangent for normal mapping
    Vec3 bitangent;          // Bitangent (normal × tangent)
    Vec4 color;              // Per-vertex color
    u8 boneIndices[4];       // Skeletal animation bone indices
    f32 boneWeights[4];      // Skeletal animation bone weights
};
\end{lstlisting}

The \texttt{normal} vector is computed as the normalized cross product of two edges emanating from the vertex. For a smooth appearance, vertex normals are typically averaged across all faces that share the vertex. This technique, known as vertex normal averaging or Phong shading, creates the illusion of a curved surface even when the underlying geometry is faceted.

\begin{equation}
    \vec{n} = \text{normalize}\left( \sum_{i \in \text{adjacent faces}} \vec{n}_i \right)
\end{equation}

Texture coordinates (also called UV coordinates) map points on the mesh surface to corresponding points in a two-dimensional texture image. The convention uses the range $[0, 1]$ for normalized coordinates, with $(0, 0)$ at the bottom-left corner of the texture and $(1, 1)$ at the top-right.

Tangent and bitangent vectors form a local coordinate frame at each vertex, perpendicular to the surface normal. These vectors are essential for normal mapping, a technique that allows the renderer to simulate fine geometric detail without increasing polygon count. The bitangent (also called binormal) is computed as:
\begin{equation}
    \vec{b} = \vec{n} \times \vec{t}
\end{equation}
where $\vec{t}$ is the tangent vector. Together, the normal, tangent, and bitangent form the Tangent-Bitangent-Normal (TBN) matrix, which transforms normal map samples from texture space into world space.

\needspace{5\baselineskip}
\subsection{Index Buffers and Triangle Strips}

To avoid redundant vertex data, modern graphics APIs employ index buffers (also called element buffers or index arrays). An index buffer is a list of integers that reference vertices in the vertex buffer. When rendering, the GPU reads vertices in the order specified by the index buffer, allowing a single vertex definition to be reused across multiple triangles.

\begin{lstlisting}[language=C++, caption=Indexed Mesh Rendering]
std::vector<Vertex> vertices = { /* ... */ };
std::vector<u32> indices = { 0, 1, 2, 2, 1, 3, /* ... */ };

// Submit to GPU
renderDevice->uploadVertexBuffer(vertices.data(), vertices.size());
renderDevice->uploadIndexBuffer(indices.data(), indices.size());
renderDevice->drawIndexed(indices.size(), 0);
\end{lstlisting}

For meshes with high polygon counts, index buffers can reduce memory usage by up to 75\% compared to duplicate vertex data. A vertex may be referenced by up to 6 adjacent triangles in a closed mesh (on average 3 references per vertex for well-formed geometry).

Triangle strips further optimize vertex submission by leveraging the connectivity of sequential triangles. In a strip, each new vertex after the first two forms a triangle with the previous two vertices. The winding order alternates to maintain consistent face orientation.

\begin{equation}
    \text{Triangles formed: } (v_0, v_1, v_2), (v_2, v_1, v_3), (v_2, v_3, v_4), \ldots
\end{equation}

Caffeine's mesh representation supports both indexed triangle lists and triangle strips, with the latter providing superior cache locality during vertex fetch.

\needspace{6\baselineskip}
\section{Primitive Shape Generation}

For rapid prototyping and editor visualization, Caffeine provides built-in procedures for generating common primitive shapes. The \texttt{MeshFilterComponent} can reference these primitives, eliminating the need for external mesh files for basic geometric forms.

\needspace{5\baselineskip}
\subsection{Cube Generation}

A cube is the simplest three-dimensional primitive. It consists of 8 vertices and 6 faces (12 triangles). Caffeine generates a unit cube centered at the origin with dimensions $[-0.5, 0.5]$ on each axis.

\begin{lstlisting}[language=C++, caption=Unit Cube Vertex Layout]
const f32 h = 0.5f;  // Half-extent
Vec3 cubeVertices[] = {
    {-h, -h, -h},  // Back-bottom-left
    { h, -h, -h},  // Back-bottom-right
    { h,  h, -h},  // Back-top-right
    {-h,  h, -h},  // Back-top-left
    {-h, -h,  h},  // Front-bottom-left
    { h, -h,  h},  // Front-bottom-right
    { h,  h,  h},  // Front-top-right
    {-h,  h,  h}   // Front-top-left
};
\end{lstlisting}

Each face of the cube must have its normal vector pointing outward. For a face aligned with the positive Z-axis, the normal is $(0, 0, 1)$. Faces are defined in counter-clockwise order when viewed from outside the cube, following the right-hand rule convention.

\needspace{5\baselineskip}
\subsection{Sphere Generation}

Spheres are generated using the UV sphere algorithm, which divides the sphere into latitude and longitude segments. Given a sphere of radius $r$, a point on the surface at latitude $\phi$ (elevation) and longitude $\theta$ (azimuth) is:

\begin{equation}
    \vec{p} = \begin{pmatrix} r \sin(\phi) \cos(\theta) \\ r \cos(\phi) \\ r \sin(\phi) \sin(\theta) \end{pmatrix}
\end{equation}

The sphere's vertices are organized in rings. Each ring corresponds to a constant latitude, and rings are connected by triangles to form the spherical surface.

\begin{lstlisting}[language=C++, caption=UV Sphere Generation, basicstyle=\ttfamily\small]
u32 latSegments = 16, lonSegments = 32;
std::vector<Vertex> vertices;

for (u32 lat = 0; lat <= latSegments; ++lat) {
    f32 phi = 3.14159265f * lat / latSegments;
    f32 sinPhi = std::sin(phi), cosPhi = std::cos(phi);
    
    for (u32 lon = 0; lon <= lonSegments; ++lon) {
        f32 theta = 2.0f * 3.14159265f * lon / lonSegments;
        f32 sinTheta = std::sin(theta), cosTheta = std::cos(theta);
        
        Vec3 pos = {
            sinPhi * cosTheta,
            cosPhi,
            sinPhi * sinTheta
        };
        
        Vertex v;
        v.position = pos;
        v.normal = pos;  // For unit sphere, normal = position
        v.texCoord = {
            static_cast<f32>(lon) / lonSegments,
            static_cast<f32>(lat) / latSegments
        };
        vertices.push_back(v);
    }
}
\end{lstlisting}

The UV sphere algorithm naturally produces texture coordinates suitable for spherical mapping. A horizontal strip around the equator maps to the full width of a texture, while the poles converge to single points in texture space.

\needspace{5\baselineskip}
\subsection{Cylinder and Capsule Generation}

A cylinder consists of two circular caps and a cylindrical side surface. The top and bottom caps are generated as triangle fans, with center vertices at $(0, \pm h/2, 0)$ and outer vertices distributed on circles of radius $r$ at those heights.

\begin{equation}
    \text{Top cap center: } (0, h/2, 0)
\end{equation}
\begin{equation}
    \text{Bottom cap center: } (0, -h/2, 0)
\end{equation}

The cylindrical surface is generated by connecting rings of vertices at the top and bottom heights.

A capsule extends the cylinder concept by replacing the flat caps with hemisphere. It is equivalent to a cylinder with hemispherical endcaps. The capsule is useful in physics simulations as a simplified collision shape for elongated objects like limbs or weapons.

\needspace{5\baselineskip}
\subsection{Plane and Quad Generation}

A plane (or quad) is the simplest two-dimensional mesh embedded in three-dimensional space. It consists of 4 vertices and 2 triangles (a single quad rendered as two triangles with a shared diagonal).

\begin{lstlisting}[language=C++, caption=Quad Mesh Generation]
const f32 w = 1.0f, h = 1.0f;
std::vector<Vertex> planeVertices = {
    {{-w/2,  0,  h/2}, {0, 1, 0}, {0, 1}},  // Top-left
    {{ w/2,  0,  h/2}, {0, 1, 0}, {1, 1}},  // Top-right
    {{ w/2,  0, -h/2}, {0, 1, 0}, {1, 0}},  // Bottom-right
    {{-w/2,  0, -h/2}, {0, 1, 0}, {0, 0}}   // Bottom-left
};
std::vector<u32> planeIndices = {
    0, 1, 2,  // First triangle
    0, 2, 3   // Second triangle
};
\end{lstlisting}

The plane's normal vector points upward $(0, 1, 0)$. When a plane is used for a ground surface or billboard, the normal determines how lighting is applied.

\needspace{6\baselineskip}
\section{Mesh Asset Loading and Caching}

The Caffeine Engine supports multiple mesh file formats, primarily glTF/glb (GL Transmission Format) and Wavefront OBJ. These formats are parsed at asset load time and converted into the engine's internal mesh representation.

\needspace{5\baselineskip}
\subsection{glTF Format Integration}

glTF is a modern, standardized format optimized for web and game engine use. It supports:
\begin{itemize}
    \item Multi-material meshes with per-surface material assignment
    \item Skeletal animation with armature definitions
    \item Embedded or external texture references
    \item Compressed binary data (glb variant)
\end{itemize}

The Caffeine Engine uses a third-party glTF parser that produces in-memory mesh data. During import, vertex data is reorganized into the engine's standard \texttt{Vertex} structure, and index buffers are validated for correctness.

\needspace{5\baselineskip}
\subsection{Mesh Instance Caching}

Once a mesh asset is loaded, it is cached in the \texttt{AssetManager}. Subsequent requests for the same mesh return a handle to the cached data, avoiding redundant file I/O and parsing.

\begin{lstlisting}[language=C++, caption=Mesh Loading with Caching]
AssetHandle<Mesh> meshHandle = assetManager->load<Mesh>("models/character.glb");

// First call: loads from disk
// Subsequent calls with same path: return cached handle
\end{lstlisting}

The mesh cache is keyed by asset path and retains loaded meshes until the asset manager is shut down or explicitly flushed. For games with large asset libraries, this cache dramatically improves load times and frame rate stability during gameplay.

\needspace{6\baselineskip}
\section{Mesh Transformation and World Space Conversion}

Entities in the scene graph maintain local (model) space coordinates. To render a mesh, the engine transforms vertex positions from model space into world space using the entity's transformation matrix.

\needspace{5\baselineskip}
\subsection{Vertex Transformation Pipeline}

A single vertex progresses through several transformation stages:

\begin{enumerate}
    \item \textbf{Model Space}: Vertex coordinates as defined in the mesh asset.
    \item \textbf{World Space}: Coordinates after applying the entity's transformation matrix (position, rotation, scale).
    \item \textbf{View Space}: Coordinates relative to the camera, after multiplying by the camera's view matrix.
    \item \textbf{Clip Space}: Coordinates after perspective projection, ready for rasterization.
\end{enumerate}

The transformation from model space to world space is accomplished by the model matrix $M$, which encodes the entity's translation, rotation, and scale:

\begin{equation}
    \vec{p}_{world} = M \cdot \vec{p}_{model}
\end{equation}

Normal vectors require special treatment. A normal must be transformed by the inverse transpose of the model matrix to ensure perpendicularity is preserved even in the presence of non-uniform scaling:

\begin{equation}
    \vec{n}_{world} = \text{normalize}((M^{-T}) \cdot \vec{n}_{model})
\end{equation}

where $M^{-T}$ denotes the inverse transpose of the model matrix.

\needspace{5\baselineskip}
\subsection{Batching and Instancing}

For scenes with many identical meshes (e.g., trees in a forest, rocks in a rock field), the engine supports mesh instancing. Instead of submitting each instance as a separate draw call, a single draw call with an instance count renders multiple copies of the mesh. Each instance can have different transformation matrices, stored in an instance buffer.

\begin{lstlisting}[language=C++, caption=Instance Buffer Layout]
struct InstanceData {
    Mat4 modelMatrix;
    Mat4 normalMatrix;  // For normal transformation
};

// GPU buffer contains array of InstanceData
// Draw call specifies instance count
renderDevice->drawIndexedInstanced(indexCount, instanceCount, 0, 0);
\end{lstlisting}

Instancing reduces CPU-to-GPU overhead significantly. Drawing 1000 trees as 1000 instanced draw calls (1 per tree) is far more efficient than 1000 separate draw calls.

\needspace{6\baselineskip}
\section{Gizmo Visualization in the Editor}

The Caffeine Editor uses gizmos—simple geometric representations—to visualize entities and components that may not have visual renderers. Light sources, camera frustums, and emitter origins are displayed as overlaid geometry in the scene viewport.

\needspace{5\baselineskip}
\subsection{Light Gizmo Representation}

Light components are visualized as follows:

\begin{itemize}
    \item \textbf{Directional Lights}: A sun symbol consisting of a filled circle with 6 radiating rays.
    \item \textbf{Point Lights}: A filled circle at the light's center with an outline circle showing the light's falloff radius.
    \item \textbf{Spot Lights}: A cone representation with a directional axis line and a circular base indicating the light cone angle.
\end{itemize}

These gizmos are drawn as 2D ImGui overlays projected onto the viewport using the editor's camera projection. The gizmo color matches the light's color, and opacity is modulated by the light's intensity.

\begin{lstlisting}[language=C++, caption=Light Gizmo Drawing]
void drawLightGizmos(ECS::World& world, EditorContext& ctx, 
                     ImVec2 origin, ImVec2 viewportSize) {
    ImDrawList* dl = ImGui::GetWindowDrawList();
    
    // Query all point lights
    ECS::ComponentQuery q;
    q.with<ECS::LightComponent>();
    q.with<ECS::PointLightComponent>();
    q.with<ECS::Transform>();
    
    world.forEach<ECS::LightComponent, ECS::PointLightComponent, 
                  ECS::Transform>(q, 
        [&](ECS::Entity entity, ECS::LightComponent& lc, 
            ECS::PointLightComponent& pl, ECS::Transform& t) {
            // Project world position to screen
            ImVec2 screenPos = projectToScreen(t.position, origin, 
                                               viewportSize, ctx);
            
            // Create color with intensity factoring
            ImU32 color = IM_COL32(
                static_cast<u8>(lc.color.x * 255),
                static_cast<u8>(lc.color.y * 255),
                static_cast<u8>(lc.color.z * 255),
                static_cast<u8>(lc.color.w * 255 * lc.intensity)
            );
            
            // Draw filled circle at position
            dl->AddCircleFilled(screenPos, 5.0f, color, 16);
            
            // Draw radius outline
            // (Convert world radius to screen radius via projection)
            dl->AddCircle(screenPos, radiusScreen, color, 16, 1.5f);
        });
}
\end{lstlisting}

\needspace{5\baselineskip}
\subsection{Camera Frustum Visualization}

Cameras used in the scene are visualized by drawing the edges of their viewing frustum. Six planes (near, far, left, right, top, bottom) define the frustum boundary. The renderer draws lines connecting the corners of the near and far planes, forming a wireframe pyramid or rectangular prism.

\begin{equation}
    \text{Near plane corners} \to \text{Far plane corners (connected by edges)}
\end{equation}

This visualization helps level designers understand the camera's view and ensures important scene elements are not accidentally outside the culled region.

\needspace{5\baselineskip}
\subsection{Physics Collider Visualization}

Physics components are rendered as wireframe shapes overlaid on the viewport. AABB colliders appear as rectangles, circles as circular outlines, and complex shapes as wireframe approximations. This debug view is toggled via the Physics Debug button in the viewport toolbar and helps developers verify collision geometry placement.

\needspace{6\baselineskip}
\section{Normal Mapping and Tangent Space}

Normal mapping allows a texture to simulate surface detail without increasing polygon count. Instead of storing actual geometric normals, a texture stores perturbed normals in a local coordinate frame called tangent space.

\needspace{5\baselineskip}
\subsection{Tangent Space Fundamentals}

At each pixel, the shader reconstructs a local coordinate frame using the vertex's normal, tangent, and bitangent vectors. Normals sampled from the normal map are in this tangent space and must be transformed to world space before being used in lighting calculations.

The transformation from tangent space to world space uses the TBN matrix:

\begin{equation}
    \vec{n}_{world} = \begin{pmatrix} t_x & b_x & n_x \\ t_y & b_y & n_y \\ t_z & b_z & n_z \end{pmatrix} \cdot \vec{n}_{tangent}
\end{equation}

where $\vec{t}$, $\vec{b}$, and $\vec{n}$ are the tangent, bitangent, and normal vectors respectively.

\needspace{5\baselineship}
\subsection{Tangent Vector Calculation}

Tangent vectors are typically computed during mesh import. For a texture coordinate $(u, v)$ varying across a triangle, the tangent direction corresponds to the direction of increasing $u$. This is computed by solving for the vectors $\vec{t}$ and $\vec{b}$ that align with texture coordinate edges.

Given two edges of a triangle:
\begin{equation}
    \vec{e}_1 = \vec{p}_1 - \vec{p}_0, \quad \vec{e}_2 = \vec{p}_2 - \vec{p}_0
\end{equation}

and their corresponding texture coordinate differences:
\begin{equation}
    \Delta u_1 = u_1 - u_0, \quad \Delta v_1 = v_1 - v_0
\end{equation}

The tangent vector can be computed by solving the linear system. The bitangent is then derived as the cross product with the normal.

\needspace{6\baselineship}
\section{Best Practices for Mesh Authoring}

Efficient and visually correct mesh rendering requires attention to several design considerations.

\needspace{5\baselineship}
\subsection{Polygon Budget}

High polygon counts improve visual fidelity but reduce performance. Developers must balance quality with frame rate requirements. Guidelines:

\begin{itemize}
    \item \textbf{Main characters}: 15,000–50,000 polygons
    \item \textbf{Secondary characters}: 5,000–15,000 polygons
    \item \textbf{Background props}: 1,000–5,000 polygons
    \item \textbf{Interactive objects}: 2,000–10,000 polygons
\end{itemize}

These budgets assume modern hardware and a target frame rate of 60 FPS with typical scene complexity.

\needspace{5\baselineskip}
\subsection{Normal Consistency}

All vertex normals must point outward for closed meshes. Inconsistent normals cause lighting artifacts and backface culling failures. Most mesh authoring tools (Blender, Maya, 3ds Max) provide automatic normal recalculation to enforce consistency.

\needspace{5\baselineskip}
\subsection{Texture Coordinate Seams}

Textures must be unwrapped onto flat 2D space (UV mapping) with minimal distortion. Seams (boundaries between UV patches) can cause visible artifacts if not carefully placed on inconspicuous parts of the model.

\needspace{5\baselineskip}
\subsection{Skeletal Structure for Animation}

Animated models require an armature (skeleton) with bones and weight assignments. Vertices near a joint should be weighted to multiple bones to create smooth deformation. A vertex weighted 50\% to the upper arm bone and 50\% to the forearm bone will deform correctly at the elbow.

\begin{equation}
    \text{Deformed Position} = \sum_{i=1}^{n} w_i \cdot B_i \cdot \vec{p}
\end{equation}

where $w_i$ is the weight for bone $i$, $B_i$ is the bone's transformation matrix, and $\vec{p}$ is the vertex position.

\needspace{6\baselineskip}
\section{Performance Optimization Techniques}

Real-time rendering of complex scenes requires careful optimization of mesh handling.

\needspace{5\baselineskip}
\subsection{Level of Detail (LOD)}

For distant objects, high-polygon meshes are unnecessary. The engine can provide multiple mesh versions at different detail levels, automatically selecting the appropriate LOD based on distance from the camera.

\begin{lstlisting}[language=C++, caption=LOD Selection]
const f32 distance = Vec3::distance(camera.position, entity.position);
u32 lodLevel = 0;
if (distance > 50.0f) lodLevel = 1;
if (distance > 100.0f) lodLevel = 2;
if (distance > 200.0f) lodLevel = 3;

Mesh* mesh = entity.getLODMesh(lodLevel);
renderDevice->drawIndexed(mesh->indexCount, 0);
\end{lstlisting}

\needspace{5\baselineskip}
\subsection{Mesh Simplification}

For visual fidelity without excessive polygon counts, mesh simplification algorithms (e.g., quadric error metrics) can automatically reduce polygon count while preserving silhouettes and important features.

\needspace{5\baselineskip}
\subsection{Vertex Buffer Pooling}

Instead of allocating new GPU buffers for every mesh, a pool of pre-allocated buffers can be reused. This reduces GPU memory fragmentation and allocation overhead.

\needspace{6\baselineskip}
\section{Future Directions}

Emerging techniques for polygon rendering include:

\begin{itemize}
    \item \textbf{Mesh Shaders}: A newer GPU feature allowing more flexible mesh processing without the traditional vertex/tessellation/geometry pipeline.
    \item \textbf{Hardware-Accelerated Ray Tracing}: Real-time ray tracing for reflections and shadows, with dedicated GPU support.
    \item \textbf{Procedural Mesh Generation}: Runtime generation of meshes for terrain, caves, and other large-scale geometry.
    \item \textbf{Compressed Vertex Formats}: Reduced precision representations (e.g., 16-bit positions, quantized normals) to reduce memory bandwidth.
\end{itemize}

These techniques will continue to push the boundaries of visual quality and performance in real-time interactive applications.
